\section*{Chapter \ref{ch:mx:the-almgren-chriss-framework} --- The Almgren--Chriss Framework}

\begin{solution}{exo:mx:the-almgren-chriss-framework:1}
$\kappa=\sqrt{10^{-6}\times1/2.53\times10^{-7}}=\sqrt{3.95}=1.99$ per day.
\end{solution}

\begin{solution}{exo:mx:the-almgren-chriss-framework:2}
$\E=\eta X^2/T=2.53\times10^{-7}\times10^{12}=253\,000$ dollars; $\operatorname{Var}=\sigma^2X^2T/3$, a standard deviation of $10^6/\sqrt3=577\,000$ dollars
(576\,000 on one-minute intervals).
\end{solution}

\begin{solution}{exo:mx:the-almgren-chriss-framework:3}
5.9 of the 6.5 hours. On a long window $x(t)\approx Xe^{-\kappa t}$, so 95\% is sold at $\ln20/\kappa$.
\end{solution}

\begin{solution}{exo:mx:the-almgren-chriss-framework:4}
The temporary cost grows like $\kappa$ (trading faster costs more per share on every share), while the risk saved by hurrying is small once the order is
already sold early; being patient adds risk slowly because the variance is bounded by the straight line's.
\end{solution}

\begin{solution}{exo:mx:the-almgren-chriss-framework:5}
Each slice's slippage measures only the book it consumed; the impact of earlier slices that has not decayed raises the prices later slices pay, a
permanent (or slowly decaying) impact the model's $\gamma=0$ leaves out.
\end{solution}

\begin{solution}{exo:mx:the-almgren-chriss-framework:6}
$\sigma$ was estimated from one-minute changes and scaled as if prices were a random walk; the market mean-reverts (its fundamentalists pull the price
back), so the variance over thirty minutes is less than thirty times the one-minute variance.
\end{solution}

\begin{solution}{exo:mx:the-almgren-chriss-framework:7}
The mean gap is $\gamma Q/2$: $2\times(5.2-1.2)/20\,000=4.0\times10^{-4}$ ticks a share for $\kappa T=0$ and $2\times(7.4-3.6)/20\,000=3.8\times10^{-4}$ for
$\kappa T=6$: consistent, as a schedule-independent permanent term should be.
\end{solution}

\begin{solution}{exo:mx:the-almgren-chriss-framework:8}
Its trades are certain; its cost is not: it holds the position the longest, and the price moves while it does. Its risk is the largest on the frontier,
a standard deviation of USD~576\,000 for the block.
\end{solution}

\begin{solution}{pb:mx:the-almgren-chriss-framework:1}
\textbf{1.} See the definitions.
\textbf{2.} $\E=\tfrac12\gamma X^2+\eta\int v^2$, $\operatorname{Var}=\sigma^2\int x^2$.
\textbf{3.} Every share sold moves the price by $\gamma$ whenever it is sold; the total is $\tfrac12\gamma X^2$.
\textbf{4.} Euler--Lagrange: $\eta\ddot x=\lambda\sigma^2x$ with $x(0)=X$, $x(T)=0$.
\textbf{5.} At 10\% participation the square-root law gives $0.8\sigma\sqrt{0.1}$ of the price; dividing by the rate gives $\eta=2.53\times10^{-7}$ dollars per share
per share a day.
\textbf{6.} $\kappa=1.99$ per day; USD~300\,000, 470\,000 and 521\,000.
\textbf{7.} After 5.9 hours.
\textbf{8.} USD~253\,000, 576\,000 and 585\,000.
\textbf{9.} The (risk, cost) pairs of optimal schedules over all risk aversions.
\textbf{10.} Patient: 258\,000, 542\,000, 551\,000 (+5.8\%); hurried: 506\,000, 352\,000, 630\,000 (+20.9\%).
\textbf{11.} Patience: the frontier is flat around the optimum toward lower \omterm{def:mx:the-almgren-chriss-framework:urgency}{urgency} and steep toward higher.
\textbf{12.} How many dollars of expected cost the manager would pay to remove a dollar squared of variance: a statement about how much the day's price risk
matters to the fund.
\textbf{13.} 20\,000 shares over 30 minutes in 30 slices at urgencies 0, 2 and 6, eight sessions each.
\textbf{14.} Means 5.2, 5.3 and 7.4 ticks; standard deviations 8.0, 6.5 and 3.8.
\textbf{15.} $\eta=0.11$, $\sigma=0.69$; means 1.2, 1.4, 3.6 and standard deviations 16.4, 13.2, 7.6.
\textbf{16.} Accumulating impact across slices (a permanent part) and mean reversion of the price.
\textbf{17.} About $4\times10^{-4}$ ticks a share.
\textbf{18.} \emph{Named result}: with $\lambda=10^{-6}$ the optimal schedule sells 95\% in 5.9 hours, expects USD~300\,000 with a standard deviation of 470\,000;
a schedule twice too patient costs 5.8\% more certainty equivalent and one twice too hurried 20.9\% more.
\textbf{19.} The ranking of schedules is reliable; the levels depend on calibrations that the simulation shows can be off by a factor of three.
\textbf{20.} The risk that comes with each cost, and so the choice of how much cost to pay to reduce it.
\end{solution}

\begin{solution}{iq:mx:the-almgren-chriss-framework:1}
Trading fast pays temporary impact; trading slowly bears price risk. The optimum front-loads the order: $x(t)=X\sinh(\kappa(T-t))/\sinh(\kappa T)$, a
straight line when risk does not matter and an exponential decay when it matters a lot.
\iqlookfor{The trade-off; the sinh shape; limits.}
\end{solution}

\begin{solution}{iq:mx:the-almgren-chriss-framework:2}
Minimise $\int(\eta\dot x^2+\lambda\sigma^2x^2)dt$; Euler--Lagrange $\eta\ddot x=\lambda\sigma^2x$; boundary conditions give the sinh ratio.
\iqlookfor{Calculus of variations; boundary conditions.}
\end{solution}

\begin{solution}{iq:mx:the-almgren-chriss-framework:3}
$\kappa$ grows by $\sqrt2$: the schedule front-loads more, the expected cost rises and the risk falls; the client moves along the frontier toward hurry.
\iqlookfor{Direction of each change; the square root.}
\end{solution}

\begin{solution}{iq:mx:the-almgren-chriss-framework:4}
Temporary: regress each slice's slippage against the mid on its trading rate; permanent: the drift of the price over the whole execution beyond the
temporary part, or the price after completion; both on many orders, with the square-root law as a check.
\iqlookfor{Slice-level and order-level estimates; sample size.}
\end{solution}

\begin{solution}{iq:mx:the-almgren-chriss-framework:5}
Given the order and the market state, return target holdings at requested times (and a way to update them as the market changes); the child-order
layer trades toward the targets. Keep it stateless or explicit about state so schedules can be tested offline.
\iqlookfor{Targets over time; separation from placement.}
\end{solution}

\begin{solution}{iq:mx:the-almgren-chriss-framework:6}
Non-linear and transient impact, intraday volume and volatility patterns, signals, spread and fees, and adaptation to the market. For a liquid stock the
volume pattern and transient impact usually matter most.
\iqlookfor{A list; a prioritisation with a reason.}
\end{solution}
